% TEMPLATE-MILC-v3.tex - plantilla maestra UNICA de las guias MILC v3.
%
% La usan las tres familias de guias, cada una con su driver:
%   · Grados 6-11  -> scripts/build-guias-g11.py
%   · Bebras       -> scripts/build-guias-bebras.py
%   · Semillero    -> scripts/build-guias-semillero.py
%
% Antes habia tres copias casi identicas (~400 lineas iguales, 6 distintas):
% cada mejora habia que aplicarla tres veces, y en la practica no se hacia
% (el motor de recursos vivia solo en dos de ellas). Lo que varia entre
% familias esta parametrizado en 6 placeholders que cada driver llena
% (se nombran aqui SIN su sintaxis real: el builder reemplaza por texto plano,
% tambien dentro de los comentarios, y un valor multilinea rompe el preambulo):
%
%   HEADER_IZQ        encabezado izquierdo de pagina
%   HEADER_DER        encabezado derecho de pagina
%   PORTADA_ETIQUETA  rotulo sobre el numero grande (GRADO/MOMENTO/MODULO)
%   PORTADA_SERIE     linea de serie de la portada
%   PORTADA_ID        identificador de la guia en la portada
%   PORTADA_CIRCULO   el circulito de grado (°), o vacio si no aplica
%
% El resto de claves las documenta content/guias/_SCHEMA.md. El script valida
% que no queden placeholders sin reemplazar antes de escribir el archivo final.

\documentclass[11pt,letterpaper]{article}
\usepackage[margin=2.0cm]{geometry}
\usepackage[table]{xcolor}
\usepackage{fontspec}
\usepackage[spanish]{babel}
\usepackage{tikz}
% Librerías TikZ para los diagramas de las guías (bloque `recursos:`):
%  · babel  → babel en español vuelve activos caracteres como «>», y TikZ los usa
%             en sus opciones (>=stealth, ->). Sin esta librería, cualquier
%             diagrama con flechas revienta con «\language@active@arg> has an
%             extra }». Debe ir ANTES de que se dibuje nada.
%  · positioning        → sintaxis «below left=2pt and 3pt».
%  · decorations.pathreplacing → llaves ({brace}) para señalar rangos.
\usetikzlibrary{babel,positioning,decorations.pathreplacing}
\usepackage{graphicx}
% Circuitos: el trabajo de electrónica se hace hoy con micro:bit y sus sensores
% integrados (MakeCode, sin cableado), así que no se necesitan esquemáticos.
% Si algún día entran montajes con protoboard, basta descomentar esta línea:
% los diagramas .tex/.tikz declarados en `recursos:` ya se incrustan solos.
% \usepackage{circuitikz}
\usepackage[most]{tcolorbox}
\usepackage{tabularx}
\usepackage{array}
\usepackage{fancyhdr}
\usepackage{titlesec}
\usepackage[document]{ragged2e}  % bandera a la derecha en todo el cuerpo
\usepackage{enumitem}
\usepackage{microtype}

% Helvetica Neue no trae la flecha U+2192, asi que cada «→» escrito literalmente
% en el YAML se compilaba a NADA, en silencio (462 flechas en 61 guias: rutas del
% tipo «paleta → tipografia → lockup» salian rotas). Mapeamos el caracter a su
% version matematica, que si tiene glifo. Asi el autor puede seguir escribiendo
% «→» comodamente en el YAML.
\usepackage{newunicodechar}
\newunicodechar{→}{\ensuremath{\rightarrow}}
% Mismo problema con los simbolos de marca y vinetas: el «✓» de «una palomita ✓»
% salia como cuadrito vacio en el PDF de todas las guias que lo usaban.
\usepackage{pifont}
\newunicodechar{✓}{\ding{51}}
\newunicodechar{✗}{\ding{55}}
\newunicodechar{✔}{\ding{52}}
\newunicodechar{•}{\textbullet}
\newunicodechar{≥}{\ensuremath{\geq}}
\newunicodechar{≤}{\ensuremath{\leq}}

\setmainfont{Helvetica Neue}
\setsansfont{Helvetica Neue}
\setlength{\parindent}{0pt}
\setlength{\parskip}{6pt}
% Sin particion de palabras: en 6.º-7.º una palabra cortada por guion al final
% de linea («Comuni-cacion») frena la lectura mucho mas que una linea corta.
\hyphenpenalty=10000
\exhyphenpenalty=10000
\pretolerance=2000
\tolerance=3000
\setlength{\emergencystretch}{3em}
\linespread{1.10}
\renewcommand{\arraystretch}{1.25}
\setlist{leftmargin=5mm,itemsep=1mm,topsep=1mm}
\newcolumntype{Y}{>{\RaggedRight\arraybackslash}X}

\definecolor{milcMagenta}{HTML}{E6007E}
\definecolor{milcVino}{HTML}{5A0038}
\definecolor{milcTurquesa}{HTML}{0093A5}
\definecolor{milcVerde}{HTML}{58A923}
\definecolor{milcMostaza}{HTML}{E5B400}
\definecolor{milcOcre}{HTML}{B86600}
\definecolor{milcNegro}{HTML}{241816}
\definecolor{milcGris}{HTML}{F7F7F4}
\definecolor{milcLinea}{HTML}{E8E2DD}
\definecolor{milcGradePrimary}{HTML}{7C3AED}
\definecolor{milcGradeDark}{HTML}{4C1D95}
\definecolor{milcGradeSoft}{HTML}{EDE9FE}
\definecolor{milcGradeAccent}{HTML}{CFE7FF}
\definecolor{milcGradeText}{HTML}{FFFFFF}
\color{milcNegro}

\pagestyle{fancy}
\fancyhf{}
\lhead{\small Tecnología e Informática · Grado 9}
\rhead{\small Guía 18 · MILC v3}
\cfoot{\small\thepage}
\renewcommand{\headrulewidth}{0pt}

\newtcolorbox{titlebox}[1]{
  enhanced,colback=#1,colframe=#1,arc=7mm,boxrule=0pt,
  left=7mm,right=7mm,top=3mm,bottom=3mm,
  before skip=8pt,after skip=8pt,
  fontupper=\sffamily\bfseries\Large\color{white}
}

\newtcolorbox{softbox}[3]{
  enhanced,colback=#2,colframe=#1,borderline west={4pt}{0pt}{#1},
  arc=2mm,boxrule=.4pt,left=4mm,right=4mm,top=3mm,bottom=3mm,
  before skip=6pt,after skip=8pt,title={#3},coltitle=white,
  colbacktitle=#1,fonttitle=\sffamily\bfseries,fontupper=\RaggedRight,
  attach boxed title to top left={xshift=3mm,yshift=-2mm},
  boxed title style={arc=2mm, boxrule=0pt}
}

\newtcolorbox{drawbox}[2]{
  enhanced,colback=white,colframe=#1,arc=4mm,boxrule=1pt,
  left=5mm,right=5mm,top=4mm,bottom=4mm,before skip=8pt,after skip=8pt,
  title={#2},coltitle=white,colbacktitle=#1,fonttitle=\sffamily\bfseries,
  fontupper=\RaggedRight,
  attach boxed title to top left={xshift=4mm,yshift=-2mm},
  boxed title style={arc=3mm, boxrule=0pt}
}

\newtcolorbox{infoband}[1]{
  enhanced,colback=#1!8,colframe=#1!50,arc=2mm,boxrule=.5pt,
  left=4mm,right=4mm,top=2mm,bottom=2mm,before skip=4pt,after skip=4pt,
  fontupper=\small\RaggedRight
}

% Puente narrativo entre secciones (contrato editorial Conéctate).
% Da continuidad: "Acabas de X. Ahora vas a Y porque Z."
% Visualmente: banda izquierda en color del grado, texto en itálica pequeña.
\newtcolorbox{puentebox}{
  enhanced,colback=milcGradeSoft,colframe=milcGradeDark,
  borderline west={3pt}{0pt}{milcGradeDark},
  arc=0mm,boxrule=0pt,left=5mm,right=4mm,top=2.5mm,bottom=2.5mm,
  before skip=4pt,after skip=8pt,
  fontupper=\itshape\small\color{milcNegro}\RaggedRight
}
% La flecha va en modo matematico a proposito: Helvetica Neue NO tiene el glifo
% U+2192, asi que \textrightarrow se compilaba a NADA («Missing character: There
% is no → in font Helvetica Neue») y el puente salia sin flecha. En modo
% matematico la flecha viene de la fuente matematica y si se dibuja.
\newcommand{\puente}[1]{%
  \begin{puentebox}%
    \textcolor{milcGradeDark}{$\rightarrow$}\hspace{2mm}#1%
  \end{puentebox}%
}

\newcommand{\linead}[1]{\noindent\color{milcLinea}\rule{#1}{0.5pt}\color{milcNegro}}
\newcommand{\respuesta}[1]{\par\vspace{2mm}\foreach \n in {1,...,#1}{\noindent\color{milcLinea}\rule{\linewidth}{0.5pt}\par\vspace{5mm}}\color{milcNegro}}
\newcommand{\checkbox}{\textcolor{milcGradeDark}{\fbox{\phantom{x}}}\hspace{5pt}}

\newcommand{\phasecard}[4]{
\begin{tcolorbox}[enhanced,colback=#3,colframe=#2,arc=3mm,boxrule=.5pt,height=3.05cm,valign=center,halign=center,left=1.5mm,right=1.5mm,top=2mm,bottom=2mm]
{\sffamily\bfseries\fontsize{22}{24}\selectfont\textcolor{#2}{#1}}\par
{\sffamily\bfseries\small #4}
\end{tcolorbox}
}

% ── Piezas del rediseno 2026 (legibilidad 6.º-11.º) ───────────────────
% Banda de estacion: reemplaza el titlebox macizo. Titulo a la izquierda,
% dato practico (tiempo/modalidad) a la derecha, en una sola linea.
\newcommand{\milcestacion}[2]{%
\begin{tcolorbox}[enhanced,colback=#1,colframe=#1,arc=2mm,boxrule=0pt,
  left=5mm,right=5mm,top=2.6mm,bottom=2.6mm,before skip=11pt,after skip=7pt]
{\sffamily\bfseries\fontsize{15}{18}\selectfont\color{white}#2}
\end{tcolorbox}}

% Tarjeta de paso numerada: sustituye las tablas «Pilar / Aplicacion», que
% eran formato de consulta docente, no de estudiante. El numero va en un
% circulo sobre el borde para que la secuencia se lea de un vistazo.
\newcommand{\milcpaso}[3]{%
\begin{tcolorbox}[enhanced,colback=white,colframe=milcLinea,arc=1.5mm,boxrule=.9pt,
  left=14mm,right=4mm,top=3mm,bottom=3mm,before skip=4pt,after skip=4pt,
  fontupper=\RaggedRight,
  overlay={\node[anchor=north west,circle,fill=#2,text=white,inner sep=0pt,
    minimum size=7.4mm,font=\sffamily\bfseries\small]
    at ([xshift=3.6mm,yshift=-3.2mm]frame.north west) {#1};}]
#3
\end{tcolorbox}}

% Casilla marcable de tamano util con lapiz (la anterior era \fbox{\phantom{x}},
% de alto variable segun el texto que la seguia).
\newcommand{\milcmarca}{\framebox[3.8mm]{\rule{0pt}{3.4mm}}}

% Fila marcable de lista suelta (checks de la estacion 1).
\newcommand{\milccheckrow}[1]{%
\par\vspace{1.8mm}\noindent
\makebox[6.4mm][l]{\raisebox{-.55mm}{\textcolor{milcNegro}{\milcmarca}}}%
\parbox[t]{\dimexpr\linewidth-6.4mm\relax}{\RaggedRight #1}\par}

% Celda marcable dentro de la rubrica (tabla de autoevaluacion). Misma
% sangria francesa, pero pensada para vivir dentro de una columna X.
\newcommand{\milcopcion}[1]{%
\makebox[5.2mm][l]{\raisebox{-.55mm}{\textcolor{milcNegro}{\milcmarca}}}%
\parbox[t]{\dimexpr\linewidth-5.2mm\relax}{\RaggedRight #1}}

% ── Recursos visuales de la guía ──────────────────────────────────────
% Declarados en el bloque `recursos:` del YAML y emitidos por el builder.
% \guiaFigura   → imagen rasterizada o PDF (foto, carta celeste, montaje).
% \guiaDiagrama → fuente TikZ/circuitikz (.tex) incrustada como vector:
%                 es la vía para circuitos y esquemáticos editables.
% #1 = ruta relativa al .tex · #2 = pie (puede ir vacío)
\newcommand{\guiaFigura}[2]{%
\begin{center}
\includegraphics[width=0.88\linewidth,height=0.38\textheight,keepaspectratio]{#1}\\[1.5mm]
{\footnotesize\itshape\textcolor{milcNegro!75}{#2}}
\end{center}
}

\newcommand{\guiaDiagrama}[2]{%
\begin{center}
\input{#1}\\[1.5mm]
{\footnotesize\itshape\textcolor{milcNegro!75}{#2}}
\end{center}
}

\begin{document}
\thispagestyle{empty}

% ── Portada ───────────────────────────────────────────────────────────
% Antes: un rectangulo de color a pagina completa. Se veia bien en pantalla,
% pero en la fotocopiadora del colegio se comia el toner y salia gris sucio.
% Ahora la banda ocupa el tercio superior y el resto va en blanco.
\begin{tikzpicture}[remember picture,overlay]
  \path[fill=milcGradePrimary]
    (current page.north west) rectangle ([yshift=-10.6cm]current page.north east);

  \node[anchor=north west,text=milcGradeText] at ([xshift=2.0cm,yshift=-1.55cm]current page.north west)
    {\sffamily\bfseries\fontsize{12}{14}\selectfont GRADO};
  \node[anchor=north west,text=milcGradeText,opacity=.85] at ([xshift=2.0cm,yshift=-2.35cm]current page.north west)
    {\sffamily\bfseries\fontsize{8.5}{10}\selectfont SERIE GUÍAS · TECNOLOGÍA E INFORMÁTICA};
  \node[anchor=north east,text=milcGradeText,opacity=.85,align=right] at ([xshift=-2.0cm,yshift=-1.55cm]current page.north east)
    {\sffamily\bfseries\fontsize{9}{12}\selectfont I.E. Sor María Juliana\\Cartago, Valle del Cauca};

  \node[anchor=west,text=milcGradeText] (gradonum) at ([xshift=1.85cm,yshift=-7.1cm]current page.north west)
    {\sffamily\bfseries\fontsize{112}{112}\selectfont 9};
  % El circulito de grado (°) solo aplica a las guias de grado («11°»); en
  % Bebras y Semillero el numero grande es un momento/modulo, no un grado.
  % Se posiciona relativo al nodo `gradonum`, no en coordenadas absolutas.
  \draw[milcGradeText,line width=1.6pt]
    ([xshift=2.0mm,yshift=-4.5mm]gradonum.north east) circle (.15cm);

  \node[anchor=west,text=milcGradeText,text width=11.4cm,align=left] at ([xshift=1.5cm,yshift=0.5cm]gradonum.east)
    {
     {\sffamily\bfseries\fontsize{9}{11}\selectfont\MakeUppercase{Período 2 · Diseño editorial digital}}\\[3mm]
     {\sffamily\bfseries\fontsize{21}{25}\selectfont\setlength{\baselineskip}{27pt}Edición y corrección ---\\[4pt]el ojo crítico}\\[3.5mm]
     {\sffamily\bfseries\fontsize{15}{18}\selectfont Guía 18-9-TIC}
    };
\end{tikzpicture}

\vspace*{9.4cm}
\vfill

{\sffamily\bfseries\fontsize{8}{10}\selectfont\textcolor{milcGradeDark}{LO QUE TE LLEVAS HOY}}

\begin{tcolorbox}[enhanced,colback=white,colframe=milcNegro,arc=2mm,boxrule=1.6pt,
  left=5mm,right=5mm,top=4mm,bottom=4mm,before skip=3pt,after skip=12pt,
  fontupper=\RaggedRight\fontsize{11}{15.5}\selectfont]
Bitácora de 10 errores corregidos sobre las propias páginas de la revista (5 de texto + 5 de diseño), cada uno con foto/captura del error, descripción y corrección aplicada
\end{tcolorbox}

\begin{tcolorbox}[enhanced,colback=milcGris,colframe=milcGris,arc=2mm,boxrule=0pt,
  left=5mm,right=5mm,top=3.5mm,bottom=3.5mm,after skip=12pt]
{\sffamily\bfseries\fontsize{8}{10}\selectfont\textcolor{milcNegro!55}{ESTA GUÍA ES DE}}\\[3mm]
\begin{tabularx}{\linewidth}{X p{3.2cm} p{3.2cm}}
\linead{\linewidth} & \linead{3.0cm} & \linead{3.0cm}\\[-1mm]
{\fontsize{8.5}{10}\selectfont\textcolor{milcNegro!55}{Nombre completo}} &
{\fontsize{8.5}{10}\selectfont\textcolor{milcNegro!55}{Grupo}} &
{\fontsize{8.5}{10}\selectfont\textcolor{milcNegro!55}{Fecha}}\\
\end{tabularx}
\end{tcolorbox}

\vfill

\noindent\textcolor{milcLinea}{\rule{\linewidth}{1pt}}\\[2mm]
{\sffamily\bfseries\fontsize{9.5}{12}\selectfont Autor: Dr. Álvaro Cárdenas Orozco}\\[1mm]
{\sffamily\fontsize{8.5}{11}\selectfont\textcolor{milcNegro!55}{Metodología MILC · apertura ancestral · pensamiento computacional · triángulo Dussel–estoicismo–Floridi}}

\newpage

\section*{Apertura: Edición y corrección --- el ojo crítico}

\begin{softbox}{milcGradeDark}{milcGradeSoft}{Saber ancestral}
Los \textbf{orfebres Quimbaya} del norte del Valle del Cauca tenían un protocolo de descarte despiadado: cada nariguera, pectoral o figura fundida pasaba por inspección visual y táctil; las que tenían burbujas internas, grietas finas o asimetrías mínimas se fundían de nuevo. No publicaban piezas dudosas. Siglos después, ese mismo \textbf{oficio del corrector} existía mucho antes de Word. En la tipografía colombiana del siglo XX --- las imprentas de Cali, Cartago, Bogotá, Medellín y el Eje Cafetero --- el corrector era \emph{casi tan importante como el diagramador}. Su trabajo: leer con un ojo distinto al del autor, encontrar lo que el autor ya no ve. Tildes faltantes, dedos cambiados, mayúsculas mal puestas, palabras repetidas, números mal escritos, espacios dobles. El corrector salvaba revistas, periódicos y libros enteros antes de imprimir. Hoy ese oficio sobrevive en la cultura editorial profesional aunque las máquinas lo asistan. Word te marca errores pero no los corrige por ti --- alguien tiene que leer con criterio.
\end{softbox}

\begin{softbox}{milcTurquesa}{milcGris}{Saber contemporáneo}
La \textbf{edición y corrección} de un proyecto editorial digital tiene dos capas: (a) \textbf{corrección de texto} (ortografía, gramática, puntuación, coherencia, claridad); (b) \textbf{corrección de diseño} (alineación, espacios, jerarquía, color, consistencia). Las dos capas se trabajan distinto. Una revista profesional dedica al menos \emph{30\% del tiempo} a edición y corrección antes de publicar. Saltarse esa fase es la diferencia entre un trabajo amateur y uno profesional, no importa cuán bueno fuera el primer borrador.
\end{softbox}

\begin{softbox}{milcMostaza}{milcGris}{Pregunta-puente}
\textit{¿Cuándo fue la última vez que volviste a leer un texto tuyo después de unas horas? ¿Qué encontraste que no habías visto al escribirlo?}
\end{softbox}

\begin{softbox}{milcVerde}{milcGris}{Saber y saber hacer}
Al terminar podrás: (1) \textbf{identificar} 5 tipos de errores comunes de texto y 5 de diseño; (2) \textbf{analizar} tu propia revista con ojo crítico distinto al del autor; (3) \textbf{evaluar} cada error según su impacto en el lector; (4) \textbf{aplicar} 10 correcciones documentadas sobre tu propio trabajo.
\end{softbox}



\small
\begin{tabularx}{\textwidth}{>{\bfseries}p{3.0cm}Y}
\rowcolor{milcGradePrimary!12}Autor & Dr. Álvaro Cárdenas Orozco\\
DBA & Producción de piezas editoriales digitales con criterio estético y ético (MEN, T\&I)\\
Referentes & Tipografía colombiana · Diseño centrado en accesibilidad · Floridi\\
Producto final & Bitácora de 10 errores corregidos sobre las propias páginas de la revista (5 de texto + 5 de diseño), cada uno con foto/captura del error, descripción y corrección aplicada\\
\end{tabularx}
\normalsize

\puente{Antes de empezar, mira el viaje completo. Las sesiones 1-7 te dieron criterio editorial, sistema visual y accesibilidad. Hoy aplicas el ojo crítico sobre tu propio trabajo --- es la sesión que diferencia amateur de profesional.}

\milcestacion{milcGradeDark}{Ruta MILC: así vamos a aprender}
\begin{tabularx}{\textwidth}{>{\centering\arraybackslash}X>{\centering\arraybackslash}X>{\centering\arraybackslash}X>{\centering\arraybackslash}X}
\phasecard{1}{milcVerde}{milcGris}{Escucha\\[-1mm]\small Observo, escucho y pregunto.} &
\phasecard{2}{milcTurquesa}{milcGris}{Entender\\[-1mm]\small Sistematizo y ordeno.} &
\phasecard{3}{milcMagenta}{milcGris}{Hacer\\[-1mm]\small Construyo y aplico.} &
\phasecard{4}{milcVino}{milcGris}{Cierre\\[-1mm]\small Evalúo y reflexiono.}
\end{tabularx}


\puente{Antes de corregir, vamos a leer 3 publicaciones reales con ojo de corrector. Verás errores en revistas serias, periódicos, posts virales. Eso te entrena el ojo para encontrar los tuyos.}

\milcestacion{milcVerde}{Estación 1 de 3 · Escucha}

\begin{softbox}{milcVerde}{milcGris}{Escena para escuchar}
\textbf{Actividad 1 · IDENTIFICA --- Errores en publicaciones reales} (15 min · individual). Busca \textbf{3 publicaciones reales} con buenos diseños generales (un periódico nacional, una revista cultural conocida, un post viral en redes con muchos likes). Con ojo de corrector, encuentra \textbf{al menos UN error} en cada una: puede ser ortográfico, de redacción confusa, de alineación incorrecta, de inconsistencia de estilo. Estos errores demuestran que incluso publicaciones serias necesitan corrección.
\end{softbox}

{\sffamily\bfseries\fontsize{8}{10}\selectfont\textcolor{milcNegro!55}{MARCA CUANDO LO HAYAS HECHO}}
\milccheckrow{Encontré al menos 1 error en un periódico/revista profesional.}
\milccheckrow{Encontré al menos 1 error en un post viral de redes.}
\milccheckrow{Reconocí que la corrección no es lujo --- es oficio incumplido por muchos.}

\begin{infoband}{milcVerde}


\textbf{Lo que va al cuaderno (Actividad 1 · IDENTIFICA).} Título: «Actividad 1 --- 3 errores en publicaciones reales». \textbf{Formato:} tabla de 3 columnas (Publicación~|~Tipo de error~|~Cómo lo corregirías), 3 filas. \textbf{Sabes que terminaste cuando} encontraste errores reales y entiendes que incluso publicaciones grandes los tienen --- por eso la corrección importa.
\end{infoband}

\puente{Ya viste errores en publicaciones reales. Ahora vamos a entender las \textbf{2 capas de corrección} (texto + diseño) y los 10 errores más comunes (5 de cada). Tener una lista clara hace la corrección sistemática, no caprichosa.}

\milcestacion{milcTurquesa}{Estación 2 de 3 · Entender}

\begin{softbox}{milcTurquesa}{milcGris}{Lo que vas a entender}
Toda revisión editorial tiene 2 capas con sus listas de errores comunes. Vamos a descomponerlas con los pilares del pensamiento computacional para que sepas qué buscar.
\end{softbox}

\milcpaso{1}{milcTurquesa}{La corrección se descompone en 2 capas con 5 errores típicos cada una. \textbf{Capa texto}: (a) tildes faltantes o sobrantes, (b) palabras repetidas en frases cercanas, (c) puntuación incorrecta (comas, puntos), (d) números o fechas mal escritos, (e) frases ambiguas o muy largas. \textbf{Capa diseño}: (a) alineación inconsistente, (b) espacios dobles entre palabras, (c) tipografía mezclada (más de tu par), (d) color sin función clara, (e) imagen pixelada o de baja calidad.}
\milcpaso{2}{milcTurquesa}{Toda corrección sigue el patrón \textbf{``leer con ojo del lector, no del autor''}: cuando tú escribes, tu cerebro completa lo que falta. Cuando otro lee, no completa. Por eso necesitas distancia (tiempo o persona) antes de corregir. Leer en voz alta es el truco más antiguo y efectivo.}
\milcpaso{3}{milcTurquesa}{Lo esencial de la corrección cabe en una pregunta: \emph{¿este error confunde al lector o solo me molesta a mí?}. Si confunde, corrige. Si solo es preferencia personal, suelta. La energía es limitada.}
\milcpaso{4}{milcTurquesa}{Pasos para corregir cualquier pieza editorial: (1) deja descansar el trabajo al menos 1 hora (mejor 1 día); (2) léelo en voz alta una vez completo; (3) revisa capa texto buscando los 5 errores comunes; (4) revisa capa diseño buscando los 5 errores comunes; (5) documenta cada corrección con antes/después; (6) lee de nuevo después de corregir.}

\begin{drawbox}{milcTurquesa}{Anatomía de una corrección documentada}
\textbf{Error detectado:} qué tipo (texto/diseño) y qué específicamente. \\[1mm] \textbf{Captura ANTES:} foto o pantallazo del error. \\[1mm] \textbf{Justificación:} por qué es error (no preferencia). \\[1mm] \textbf{Corrección aplicada:} qué cambiaste exactamente. \\[1mm] \textbf{Captura DESPUÉS:} resultado corregido. \\[1mm] \textbf{Impacto en el lector:} qué se gana con la corrección.
\end{drawbox}

\begin{softbox}{milcMostaza}{milcGris}{Errores comunes y su corrección}
(1) Confundir corrección con reescritura: el corrector no rehace el texto, solo lo limpia. (2) Solo confiar en Word: muchos errores los pasa de largo. (3) Corregir sin descanso: el cansancio te hace pasar errores obvios. (4) Corregir solo lo que te molesta a ti: el criterio es el lector. (5) Saltarse el paso de leer en voz alta: es la prueba más efectiva.
\end{softbox}

\begin{infoband}{milcTurquesa}


\textbf{Lo que va al cuaderno (Actividad 2 · ANALIZA).} Título: «Actividad 2 --- 10 errores comunes (5 de texto + 5 de diseño)». \textbf{Formato:} 2 columnas (Texto~|~Diseño) con 5 errores cada una, cada uno con ejemplo concreto inventado o real. \textbf{Sabes que terminaste cuando} reconoces los 10 errores en cualquier publicación que mires.
\end{infoband}

\puente{Tienes el método. Toca aplicar: revisar tu propia revista con el ojo del corrector y documentar 10 correcciones (5 de texto + 5 de diseño) con foto/captura del antes y el después.}

\milcestacion{milcMagenta}{Estación 3 de 3 · Hacer}

\begin{softbox}{milcMagenta}{milcGris}{Cómo vas a construir tu producto}
\textbf{Actividad 3 · EVALÚA --- 10 correcciones sobre tu propia revista} (30 min · individual). Hora de aplicar el ojo crítico a tu propio trabajo. Vas a tomar tus spreads (sesiones 5 y 6) y portada (sesión 6) y encontrar \textbf{10 errores} para corregir: 5 de texto + 5 de diseño. Cada uno documentado con antes/después.
\end{softbox}

\milcpaso{1}{milcMagenta}{Tu revisión se descompone en 4 momentos: (1) descansar (cierra el archivo al menos 1 hora antes de empezar a corregir), (2) leer en voz alta el spread completo, (3) buscar los 5 errores de texto con lista de control, (4) buscar los 5 errores de diseño con lista de control.}
\milcpaso{2}{milcMagenta}{Sigue el patrón \textbf{``humilde antes que defensivo''}: tu trabajo seguramente tiene los errores comunes. No te lo tomes personal --- todos los tenemos. La diferencia es si los corriges o los ignoras. Humildad ante el propio trabajo es señal de oficio.}
\milcpaso{3}{milcMagenta}{Cada corrección expresa una decisión por el lector. Si dudas si corregir, pregúntate: ``¿esto confunde al lector?''. Si sí, corrige. Si no, sigue.}
\milcpaso{4}{milcMagenta}{Algoritmo: (1) toma descanso de tus piezas; (2) lee en voz alta; (3) marca con lápiz cada error que encuentres; (4) llega a 10 mínimo (si encuentras menos, sigue buscando); (5) documenta cada uno con captura antes/después; (6) anota qué tipo de error y por qué es error; (7) revisa el resultado final.}

\begin{drawbox}{milcMagenta}{Checklist antes de cerrar la bitácora de corrección}
\checkbox Descansé mínimo 1 hora antes de corregir. \\[1mm] \checkbox Leí mi trabajo en voz alta al menos una vez completo. \\[1mm] \checkbox Encontré y corregí mínimo 5 errores de texto. \\[1mm] \checkbox Encontré y corregí mínimo 5 errores de diseño. \\[1mm] \checkbox Cada corrección tiene captura ANTES y DESPUÉS. \\[1mm] \checkbox Cada corrección tiene justificación de por qué era error.
\end{drawbox}

\begin{softbox}{milcMagenta}{milcGris}{Plantilla de trabajo}
\textbf{Corrección 1.} \textit{Tipo (texto/diseño) + descripción del error + por qué es error.} \\[1mm] \textbf{Corrección 2-10.} \textit{Lo mismo, con capturas antes/después.} \\[1mm] \textbf{Reflexión final.} \textit{¿Qué tipo de error fue más frecuente en mi trabajo y qué me dice eso?}
\end{softbox}

\begin{infoband}{milcMagenta}


\textbf{Lo que va al cuaderno (Actividad 3 · EVALÚA).} Título: «Actividad 3 --- Bitácora de 10 correcciones». \textbf{Formato:} 10 fichas mini con captura antes + después + tipo de error + justificación + reflexión final de 3 renglones. \textbf{Extensión:} 1.5 páginas de cuaderno. \textbf{Sabes que terminaste cuando} la bitácora demuestra disciplina, no improvisación.
\end{infoband}

\puente{Lo que sigue es el resumen de qué entregas hoy: bitácora de 10 correcciones. El criterio principal: que cada error sea específico (no ``ortografía general'') y que cada corrección esté documentada.}

\milcestacion{milcMostaza}{Producto de la sesión}
\textbf{Bitácora de 10 correcciones documentadas (5 texto + 5 diseño) sobre la propia revista}.
\begin{softbox}{milcMostaza}{milcGris}{Criterios de calidad}
(1) 10 correcciones documentadas (mínimo 5 de texto + 5 de diseño). (2) Cada una con captura ANTES y DESPUÉS. (3) Cada una con tipo de error y justificación. (4) Reflexión final identifica patrón en mis errores. (5) La bitácora demuestra que descansé y leí en voz alta antes.
\end{softbox}

\puente{Antes de cerrar, mira el oficio del corrector desde las cinco dimensiones humanas. La corrección parece técnica pero es profundamente ética: cuidado por el lector, disciplina personal, respeto por la palabra.}

\milcestacion{milcVino}{Evaluación liberadora · cinco dimensiones}

\begin{softbox}{milcVino}{milcGris}{Sentido de la evaluación}
Evaluar no es solo revisar una nota. Es reconocer qué aprendí, qué cambió en mí, qué emociones manejé, cómo me relaciono con la ciudad y la comunidad, y qué le dejaré a quien viene detrás.
\end{softbox}

\textbf{Mi reflexión liberadora en cinco dimensiones.} Completa cada frase iniciada:

\small
\begin{tabularx}{\textwidth}{>{\bfseries\RaggedRight\arraybackslash}p{3.4cm}Y>{\RaggedRight\arraybackslash}p{4.6cm}}
\rowcolor{milcVino}\color{white}Dimensión & \color{white}\bfseries Mi respuesta (frame starter) & \color{white}\bfseries Algo que descubrí\\
1. Personal & Lo que más cambió en mí fue \linead{6.5cm} & \linead{4.2cm}\\
2. Emocional & La emoción más fuerte fue \linead{2.5cm} y la regulé \linead{3cm} & \linead{4.2cm}\\
3. Ciudadana & En democracia esto importa porque \linead{6cm} & \linead{4.2cm}\\
4. Local (Cartago) & En mi barrio sería útil para \linead{6.5cm} & \linead{4.2cm}\\
5. Intergeneracional & A mi abuela/abuelo le diría \linead{6.5cm} & \linead{4.2cm}\\
\end{tabularx}
\normalsize

\begin{infoband}{milcVino}
\textbf{Para la próxima sustentación.} Vuelve a esta tabla en seis meses. Verás que las respuestas cambian — no porque lo aprendido cambie, sino porque tú cambias. La evaluación liberadora es una conversación contigo a través del tiempo.
\end{infoband}


\puente{Y para terminar, tres voces sobre el ojo crítico. Dussel sobre la lengua de la corrección, Marco Aurelio sobre la disciplina personal, Floridi sobre la calidad informacional.}

\milcestacion{milcGradeDark}{Triángulo de pensamiento · cierre filosófico}

\begin{softbox}{milcVino}{milcGris}{Dussel · lente del nosotros}
\textit{Toda corrección lingüística contiene una decisión política sobre qué lengua se considera legítima.}

{\footnotesize\itshape\color{milcNegro!70}--- formulación de esta guía, al modo de Enrique Dussel. No es una cita textual.}

La ``corrección'' de un texto en español de Colombia no es neutra: a veces ``corrige'' regionalismos del Valle, del Pacífico, del Caribe que son español igualmente válido. El corrector consciente respeta la voz del autor cuando es decisión deliberada y solo corrige errores genuinos. Aprende a distinguir entre error y variedad lingüística legítima --- esa distinción es ética profunda.

\textbf{¿Cuando corrijo, estoy quitando errores o también borrando voces regionales? ¿Qué se preserva y qué se uniforma?}
\end{softbox}
\linead{\linewidth}\\[1mm]
\linead{\linewidth}

\begin{softbox}{milcMostaza}{milcGris}{Estoicismo · Marco Aurelio · cuidado interior}
\textit{El cuidado paciente de lo que parece pequeño es la mayor virtud del oficio.}

{\footnotesize\itshape\color{milcNegro!70}--- formulación de esta guía, al modo de Marco Aurelio. No es una cita textual.}

Buscar una tilde faltante o un espacio doble parece tarea menor. Pero el cuidado paciente por lo pequeño es la diferencia entre un trabajo amateur y uno profesional. La sabiduría estoica del corrector: cada detalle es oportunidad de virtud. Quien no cuida los detalles, no cuida el conjunto.

\textbf{¿Me canso de buscar errores pequeños o veo en cada uno una oportunidad de oficio? ¿Qué cambia si los tomo como práctica espiritual y no como tarea?}
\end{softbox}
\linead{\linewidth}\\[1mm]
\linead{\linewidth}

\begin{softbox}{milcTurquesa}{milcGris}{Floridi · lente de la infoesfera}
\textit{La calidad informacional en la infoesfera depende de millones de actos de corrección invisibles.}

{\footnotesize\itshape\color{milcNegro!70}--- formulación de esta guía, al modo de Luciano Floridi. No es una cita textual.}

Cada texto corregido contribuye a una infoesfera más legible. Cada error no corregido es ruido que se multiplica. Floridi vería el oficio del corrector como mantenimiento ético de bien común informacional. Tu corrección sobre tu propia revista escolar es entrenamiento para ese cuidado mayor.

\textbf{¿Estoy contribuyendo a una infoesfera más legible o le sumo ruido cuando publico sin corregir? ¿Qué cambia si me trato como corrector profesional aunque sea estudiante?}
\end{softbox}
\linead{\linewidth}\\[1mm]
\linead{\linewidth}

\begin{infoband}{milcNegro}
\textbf{Nota para el docente.} Las tres formulaciones son del autor de la guía, no citas textuales de los pensadores: recogen su lente para pensar el tema, no sus palabras. Si vas a citarlos en clase, remite a la obra original.
\end{infoband}

\milcestacion{milcVino}{Mi autoevaluación · marca una casilla por fila}

{\small
\setlength{\extrarowheight}{2.2pt}
\begin{tabularx}{\textwidth}{>{\RaggedRight\arraybackslash\bfseries}p{3.0cm}YYY}
\rowcolor{milcVino}\color{white}\bfseries Criterio & \color{white}\bfseries Lo logré & \color{white}\bfseries En proceso & \color{white}\bfseries Necesito apoyo\\[.6mm]
Comprensión del tema & \milcopcion{Explico las ideas con mis palabras.} & \milcopcion{Reconozco las ideas, falta precisión.} & \milcopcion{Necesito repasar lo básico.}\\[.8mm]
\rowcolor{milcGris}Pensamiento computacional & \milcopcion{Usé los pilares al armar mi producto.} & \milcopcion{Usé algunos pilares.} & \milcopcion{Necesito releer la estación 2.}\\[.8mm]
Aplicación y producto & \milcopcion{Mi producto cumple los criterios.} & \milcopcion{Está armado, con ajustes pendientes.} & \milcopcion{Aún está en construcción.}\\[.8mm]
\rowcolor{milcGris}Reflexión 5 dimensiones & \milcopcion{Respondí cada una con honestidad.} & \milcopcion{Respondí algunas con detalle.} & \milcopcion{Mis respuestas son muy generales.}\\[.8mm]
Triángulo de pensamiento & \milcopcion{Conecté las 3 voces con mi trabajo.} & \milcopcion{Conecté al menos una.} & \milcopcion{No las relaciono con mi proceso.}\\[.8mm]
\end{tabularx}}

\vspace{3mm}
\textbf{Hoy entendí que:}\\[1mm]
\linead{\linewidth}\\[1mm]
\linead{\linewidth}

\textbf{Mi compromiso para la próxima sesión es:}\\[1mm]
\linead{\linewidth}\\[1mm]
\linead{\linewidth}

\begin{softbox}{milcMostaza}{milcGris}{Cita de despedida · Marco Aurelio}
\textit{``El obstáculo en el camino se vuelve el camino. Lo que estorba al camino se vuelve camino.''}\\[1mm]
Cada error de hoy, bien observado, se convierte en pieza del camino que pisarás mañana.
\end{softbox}

\small
\begin{tabularx}{\textwidth}{>{\bfseries}p{3.6cm}Y}
\rowcolor{milcGradeDark!12}Nombre completo: & \linead{10cm}\\
Fecha: & \linead{4cm} \quad Curso/grupo: \linead{4cm}\\
Mi firma: & \linead{9cm}\\
\end{tabularx}
\normalsize

\begin{infoband}{milcGradeDark}
\textbf{Para la próxima semana.} Aplica el ojo del corrector a UN texto tuyo que vayas a publicar esta semana (post, correo, trabajo escolar). Descansa, lee en voz alta, busca los 10 errores comunes. Anota qué encontraste. Esa práctica se vuelve hábito si la repites 5 veces.
\end{infoband}

\end{document}
